\makeatletter
\def\input@path{{../../.format/}}
\makeatother

\documentclass[runningheads]{llncs}
\usepackage[T1]{fontenc}
\usepackage{graphicx}
\usepackage{amsmath}
\usepackage{amssymb}

\begin{document}

\begin{center}
\textbf{Imaginary Reflection: On Narratives and Philosophy}
\end{center}

\vspace{0.5em}

\noindent\textit{[Observation]}\\
\noindent\textbf{Step 1:} \hfill \textit{Historical Pattern}\\
\noindent$\mathcal{M} = \{T_i \mid i \in \mathbb{N}\}$\\
\textit{Mathematics consists of multiple theories that evolved through time...}

\vspace{0.8em}

\noindent\textit{[Analysis]}\\
\noindent\textbf{Step 2:} \hfill \textit{Theory Evolution}\\
\noindent$T_{i+1} = T_i \mid\!\Rightarrow_{M_i} T_{i+1}$\\
\textit{Each new theory emerges from previous ones through metalogical transduction.}

\vspace{0.8em}

\noindent\textit{[Definition]}\\
\noindent\textbf{Step 3:} \hfill \textit{Narrative Construction}\\
\noindent$\mathcal{N} = \langle T_0, T_1, \ldots, T_n \rangle$\\
\textit{A narrative is a sequence of theories connected by transduction steps.}

\vspace{0.8em}

\noindent\textit{[Insight]}\\
\noindent\textbf{Step 4:} \hfill \textit{Mathematical History}\\
\noindent$\text{Parallel Postulate} \Rightarrow \{\text{Euclidean}, \text{Hyperbolic}, \text{Elliptic}\}$\\
\textit{The undecidability of the parallel postulate led to multiple geometries...}
\noindent\textit{[Example]}

\vspace{0.8em}

\noindent\textit{[Insight]}\\
\noindent\textbf{Step 5:} \hfill \textit{Mathematical History}\\
\noindent$\text{Axiom of Choice} \Rightarrow \{\text{ZF}+\text{AC}, \text{ZF}+\neg\text{AC}\}$\\
\textit{The undecidability of the axiom of choice led to different set theories...}
\noindent\textit{[Example]}

\vspace{0.8em}

\noindent\textit{[Recognition]}\\
\noindent\textbf{Step 6:} \hfill \textit{Narrative Patterns}\\
\noindent$\forall \mathcal{N} \exists G_{\mathcal{N}} : \text{Undecidable}(G_{\mathcal{N}}, \mathcal{N})$\\
\textit{Every mathematical narrative encounters undecidable propositions.}

\vspace{0.8em}

\noindent\textit{[Categorization]}\\
\noindent\textbf{Step 7:} \hfill \textit{Response Patterns}\\
\noindent$\mathcal{R} = \{\text{Branch}, \text{Extend}, \text{Reframe}, \text{Ignore}\}$\\
\textit{Mathematicians respond to undecidability in various ways...}

\vspace{0.8em}

\noindent\textit{[Realization]}\\
\noindent\textbf{Step 8:} \hfill \textit{Philosophical Emergence}\\
\noindent$\Phi(\mathcal{N}) = \text{Cl}\left(\bigcup_{T_i \in \mathcal{N}} T_i\right)$\\
\textit{A philosophy is the closure of all theories in a narrative.}

\vspace{0.8em}

\noindent\textit{[Connection]}\\
\noindent\textbf{Step 9:} \hfill \textit{Human Mathematics}\\
\noindent$\text{Human Mathematics} \cong \text{Narrative Construction via Transduction}$\\
\textit{The way humans develop mathematics mirrors consciousness...}

\vspace{0.8em}

\noindent\textit{[Insight]}\\
\noindent\textbf{Step 10:} \hfill \textit{Personal Evolution}\\
\noindent$\text{Learning} \cong \text{Self-Narrative Extension}$\\
\textit{Personal intellectual growth follows the same pattern of narrative extension.}

\vspace{0.8em}

\noindent\textit{[Synthesis]}\\
\noindent\textbf{Step 11:} \hfill \textit{Meta-Narrative}\\
\noindent$\text{Metalogical Inference} = \{\mathcal{N} \mid \mathcal{N} \text{ follows transduction rules}\}$\\
\textit{The field of Metalogical Inference studies how narratives evolve through transduction.}

\vspace{0.8em}

\noindent\textit{[Realization]}\\
\noindent\textbf{Step 12:} \hfill \textit{Ultimate Insight}\\
\noindent$\text{Consciousness} \cong \text{Narrative Evolution}$\\
\textit{To be conscious is to maintain and extend a narrative through transduction!}

\end{document}